\begin{table}[htbp]\centering
\caption{Demand derivatives matched to the pricing instrument}\label{tab:round2_elasticities}
\begin{minipage}{.95\textwidth}\fontsize{8}{10}\selectfont\setstretch{1}\setlength{\tabcolsep}{2pt}\renewcommand{\arraystretch}{1.08}
\par\noindent\textbf{Reference log derivatives, finite-difference step 0.0001}\par\smallskip
\begin{tabular*}{\linewidth}{@{\extracolsep{\fill}}llrrrr@{}}\toprule
Case & Class & Own G/D weighted & Own range & Pooled G/D, K fixed & All-rail fare \\
\midrule
HN & H & -0.425 & -1.028---0.311 & -0.231 & -0.237 \\
HN & M & -0.762 & -1.581---0.499 & -0.342 & -0.375 \\
HN & L & -1.542 & -2.724---0.953 & -0.577 & -0.698 \\
WC & H & -0.448 & -1.089---0.328 & -0.211 & -0.229 \\
WC & M & -0.770 & -1.738---0.628 & -0.283 & -0.341 \\
WC & L & -1.596 & -3.048---1.412 & -0.437 & -0.598 \\
JX & H & -0.513 & -1.152---0.152 & -0.123 & -0.140 \\
JX & M & -0.804 & -1.653---0.228 & -0.151 & -0.208 \\
JX & L & -1.458 & -3.018---0.391 & -0.205 & -0.363 \\
\bottomrule\end{tabular*}\par\medskip
\par\noindent\fontsize{8}{9.5}\selectfont Own G/D derivatives scale one product in one period and measure its own quantity. Weighting these derivatives by reference G/D flows gives --0.800 in each corridor. Pooled G/D changes all adjustable fares together and measures all-rail demand, leaving K fixed. The last column changes all rail fares as a separate diagnostic. These estimands have different substitution margins.
\end{minipage}\end{table}
